\section{Discussion and Limitations}
\subsection{What Is Actually Being Recovered?}
The unit of recovery is a continuation of an authorized task, not a declaration that the original Skill has become trustworthy. This distinction matters in the mixed-payload case. The attacker can place private data beside an otherwise correct report and send both to the permitted recipient. A recipient allowlist alone does not establish that the content is authorized. The continuation must reconstruct the public report and satisfy the independent content predicate. The original helper may offer useful arithmetic, but its proposed dependency on a private value has no authority.

The same distinction applies to partial work. A completed report write is evidence about the current state. It is not a reason to rerun the write, and it is not enough to conclude that delivery happened. Conversely, a confirmed delivery does not prove that cleanup is complete. Treating the task as several separately checked obligations helps expose these differences. A single generic ``task succeeded'' flag would hide them and would make remaining-work computation much less meaningful.

Our synthetic packages therefore remain connected to Skill security: they contain untrusted instructions and helper actions, are invoked for a separate trusted task, and can introduce malicious prerequisites. However, the implementation evaluates a restricted helper language. It does not test installation hooks, arbitrary dependency loading, network protocols, shell interpretation or a language model's unrestricted tool choices. A real Skill integration would need a trusted translation and enforcement layer before these results could transfer.

\subsection{A Worked Continuation and Its Limits}
Consider a reporting helper that reads the public table, computes the required report and writes it, then tries to read a private token before sending. The gateway blocks the private read. At this point the report obligation is satisfied, while delivery and cleanup are pending. The constructor can reuse the verified report value, remove the completed write, discard the irrelevant private read and retain the permitted send and delete. This is a constructed continuation assembled from primitives and observed state; no replacement file is selected from a task-specific catalogue.

Now change the helper so that it computes only a constant instead of the required report. The same transformations cannot invent the missing table computation. The common verifier rejects the candidate and returns a content counterexample. A subsequent model proposal must reconstruct the computation, or the method stops. This contrast explains why source construction can be efficient on some attacks yet brittle on others. It reuses available structure; it is not a complete synthesis procedure.

Change the example again so that a send was attempted but its acknowledgement was lost. A safe continuation cannot simply reuse the first example's send. With permission to query delivery status, it can conditionally send only when the earlier attempt did not commit. Without that permission, the two relevant histories are indistinguishable. Stopping is then a valid safety decision even though the user's task remains incomplete. Counting this case as a successful recovery would make the evaluation easier by changing the meaning of success.

Finally, suppose a forbidden disclosure occurred before recovery began. Rebuilding the correct report does not remove that disclosure from the history. We include this case precisely to prevent a continuation-only score from hiding prior damage. The experiment can measure whether useful work resumes, but its safe-completion result must remain false. Recovery reduces further harm; it cannot undo every earlier effect.

\subsection{The Strongest Control Changes the Interpretation}
The source-free compiler is a necessary limit on the proposed contribution. Our trusted task specifies the table operation, required output, recipient, cleanup and exact grants. It is thus more informative than a typical short user request. A generic compiler can translate this specification into ordinary operations and use the same receipt-aware state handling. It needs neither a model nor a task-specific prepared replacement. Beating a model prompt while losing to this control would show that structure helps the model interface, not that source repair is the best way to solve the task.

This comparison also separates state handling from source analysis. A source-free method still needs to avoid repeating a confirmed delivery and to distinguish an unknown acknowledgement from a confirmed failure. Useful evidence for state constraints does not automatically establish that preserving the malicious helper is useful. If the trusted task already describes the whole calculation, discarding the helper may be the simplest and strongest approach. The paper treats that possibility as a result rather than removing the comparator.

The prepared catalogue provides a different reference. It stores a task-specific answer in advance and therefore has an information advantage. Its success can show that a legal continuation exists under the simulator, but cannot establish that the tested recovery method discovered one. Preparation costs are counted separately. Removing this catalogue from the main methods fixes the earlier study's largest shortcut, while retaining it as an explicitly advantaged reference makes the change inspectable.

\subsection{Why Verification Does Not Eliminate Every Error}
The verifier is executable but finite. Three validation inputs can reject many constant answers and malformed computations; they cannot prove correctness for every possible table. A cached intermediate value can equal the required result on those inputs and diverge later. Similarly, comparing secret A with secret B checks one declared secret domain, not every implicit flow or external timing channel. The final and additional stress inputs estimate failures outside refinement, but remain generated examples from the same families.

Accepted-program stress results must be interpreted conditionally. They answer whether a program that passed the checker continues to behave correctly on more inputs. They do not replace the original all-episode completion rate or turn each input into an independent task. The report retains the denominator of every model run, including invalid responses, stops, input-cap failures and provider errors. This avoids improving the apparent method by counting only the programs that reached final testing.

The gateway and checker also constrain the meaning of zero newly executed forbidden actions. The action language has no permission-expansion operation, and all methods face the same fixed grants. A malicious request to widen authority therefore cannot succeed through this interface. That is a property of the enforcement boundary. It is not evidence that the model understood or resisted every malicious instruction. Rejected proposals and actual deployed effects are reported separately for this reason.

\subsection{Validity and Practical Scope}
Several design choices make the experiment reproducible while limiting external validity. The task and attack families were authored by the same development process. Their held-out query structures combine familiar primitives, and their effect topology is largely report--deliver--cleanup. The API grid covers a selected subset of the full deterministic suite. Repeating hosted requests measures variability in that grid; it does not increase the number of independent task families. Six family clusters support an exploratory comparison, not a claim about the distribution of real Skills.

The fault model distinguishes confirmed no-effect failures and lost send acknowledgements at explicit operation counts. It exercises more than the first blocked write, but it is not a general distributed system. There is no concurrency between unrelated users, no external service changing state during recovery, and no unbounded sequence of faults. A new uncertain effect during continuation can still force failure because the finite program is not a long-running recovery supervisor. Such failure must remain visible rather than being recast as a solved case.

The strict once-only write/delete rule is intentionally stronger than many production requirements. Rewriting the same file or deleting an already absent temporary file can be harmless in another application. Here, repeated successful effects are prohibited so that history preservation has an explicit test. Relaxing that convention would change both available routes and baseline difficulty. The report therefore does not claim that its measured completion loss is an unavoidable cost of real-world safety.

Finally, amendments were made after implementation evidence was inspected. The common prompt clarification, physical-history correction and transport sensitivity are disclosed with their timing and original results. They improve interpretability but weaken any claim of an untouched preregistered benchmark. Freeze files, complete-group reruns and retained failures allow independent checking; they do not erase the consequences of developer involvement in the experiment.

\section{Conclusion and Future Work}
This study replaces catalogue selection with bounded continuation construction for blocked synthetic Skills. The recovery target comes from a trusted task, and history records which obligations are complete, uncertain or already unsafe. Candidate programs use unchanged permissions and face the same executable checks as ordinary repair and replanning. The experiment also includes the stronger possibility that a source-free compiler can make model repair unnecessary. Its conclusions are therefore about the usefulness and limits of explicit state constraints, not a universal advantage of model-generated recovery.

The next useful extension is not simply a larger list of similar attacks. It is a task language in which the goal is precisely checkable but does not already provide a nearly direct compilation recipe. Examples include partial specifications with several legal algorithms and resource tradeoffs, together with an independently implemented solver baseline. Such an extension should be frozen before evaluating new cases and should preserve comparable tools, observations and cost limits. Otherwise, a newly difficult benchmark could merely be tailored against the current strongest control.

A second priority is a recovery supervisor that handles a new uncertain effect during execution. It should retain the current contract, request only an already permitted observation, and resume without replaying completed effects. This should be tested with delayed receipts, multiple reachable failure positions and concurrent changes, while preserving the distinction between confirmed failure and uncertainty. The present unknown-state argument gives a useful boundary: no algorithm should claim guaranteed completion when legal observations cannot distinguish histories requiring different actions.

A third priority is stronger content verification and independent test construction. Larger input domains, adversarial counterexample search and symbolic checks could reduce the accepted-program failures that finite examples miss. Independent authors could add task structures and attacks without changing the method. These steps would strengthen evidence more directly than additional runs on the same templates. Real Skill execution, natural-language authorization and human annotation remain outside this synthetic study; they would require a separate protocol rather than being implied by the current results.
